\documentclass[instructions.tex]{subfiles}
\begin{document}

\chapter{De la haine de Dieu pour le péché.}


\primo Il n'y a que le péché qui soit l'objet de la haine de Dieu: l'enfer même n'est point haï de Dieu, parce que Dieu ne hait point ses ouvrages\footnote{Nihil odisti eorum quæ fecisti. Sap.11.}. Or, l'enfer est son ouvrage, et non pas le péché.

Autant que Dieu s'aime lui-même, autant il hait le péché; et, s'il pouvait cesser de le haïr, il se contredirait lui-même. Dieu est le souverain bien; il est donc souverainement opposé au péché, qui est le souverain mal. Dieu est le souverain ordre, la souveraine règle, la souveraine sainteté; il doit donc haïr le péché, qui est un dérèglement et un désordre souverainement opposé à sa sainteté. Il est aussi impossible que Dieu aime le péché, qu'il est impossible qu'il se haïsse lui-même. La haine que Dieu a pour le péché est donc une haine essentielle, une haine éternelle, une haine souveraine et infinie.

\secundo Je dois conclure de cette vérité, que je dois haïr le péché de toutes mes forces. C'est en vain que je me flatte d'être uni à Dieu et d'être dans sa grâce, si le péché mortel règne en moi. Le péché ne peut approcher du trône de Dieu, ni monter dans le ciel; de même Dieu ne peut habiter, par son amour, dans un cœur où règne le péché. Il est donc vrai que je suis ennemi de Dieu et haï de Dieu, si j'ai le malheur d'être en péché.

Comprenons-nous combien cet état est horrible? On a vu de grands hommes mourir de douleur, parce qu'ils étaient disgraciés de leur prince. Hélas! que craignons-nous sur la terre, si nous ne craignons une disgrâce bien plus redoutable, qui est la haine de Dieu? Être haï de Dieu est un mal plus à craindre que l'enfer même. L'enfer n'a rien de plus affreux que le malheur d'être haï de Dieu; et il ne serait plus un enfer, si l'on y était aimé de Dieu.


\section{Résolutions}

Je ne haïrai rien tant en ce monde que le péché.

\primo Et tous les jours je demanderai à Dieu qu'il veuille bien augmenter en moi cette haine si nécessaire.

\secundo Mon Dieu, faites-moi la grâce d'être prêt constamment à souffrir en cette vie, excepté le péché.

\end{document}
